
\subsection{Simulation Network}


\begin{figure}[htbp]
    \centering
    \includegraphics[width=0.85\linewidth]{network/network_original_paper.png}
    \caption{The simulated on-ramp merging area. Figure taken from \cite{hou2021rampmeter}.}
    \label{fig:network}
\end{figure}




The simulation network consists of a two-lane freeway and a single-lane on-ramp merging into the mainline.
The freeway is divided into eight segments (Seg. 1 - 8), each with a length of 100 m, except for the merging area (Seg. 6), which has an extended length of 150 m to enable smoother merging. 

The on-ramp consists of additional segments, Seg. 9 and Seg. 10, with lengths of 100 m and 75 m, respectively. 
The ramp meter traffic signal, which is regulated by the implemented control algorithms, is located between Seg. 9 and Seg. 10. 
At Seg. 6, the network temporarily expands to three lanes as this is where the on-ramp lane gets merged into the freeway, before it tapers down to two lanes again. 

\subsection{Episode Logic and Traffic Flow Generation}

To obtain statistically meaningful results, all controllers were evaluated over 20 simulation episodes (using the microscopic traffic simulator SUMO together with the TraCI interface \cite{krajzewicz2012recent})
and the reported metrics were averaged across all runs. Each episode consisted of 1000 simulation steps. Before the actual evaluation started, an additional warm-up phase of 40 
simulation steps was introduced in order to establish realistic traffic conditions before measurements took place. 

Two traffic scenarios were considered: mixed near-congested/congested traffic conditions and extremely congested traffic conditions. 
Following the assumptions made in the original study, the bottleneck capacity of the network using the default SUMO vehicle parameters was estimated to be between 3600 and 3700 veh/h. 

In the mixed traffic scenario, the freeway inflow was sampled from a Gaussian distribution $\mathcal{N}(3000, 300)$ veh/h and then clipped to the interval $[0, 3900]$ veh/h. Similarly, the on-ramp inflow was generated independently from $\mathcal{N}(900, 300)$ veh/h and clipped to $[0, 1500]$ veh/h.  
In the extremely congested traffic scenario, the freeway inflow was fixed at $3900$ veh/h, while the on-ramp inflow followed the same distribution as in the mixed traffic scenario, namely $\mathcal{N}(900, 300)$ veh/h clipped to the interval $[0, 1500]$ veh/h.

To evaluate the controllers under varying traffic conditions, the traffic inflow on both the freeway and the on-ramp was updated every 100 simulation steps according to the above distributions. 


\subsection{Emission and Energy Model}

For measuring the three metrics average vehicle fuel efficiency (mpg), average CO$_2$ emissions per vehicle (g/mi) and average NO$_x$ emissions per vehicle (mg/mi), the default emission model in SUMO \cite{krajzewicz2014sumo} was used, as specified in the original paper. 
